\documentclass[11pt]{article}
\usepackage[margin=1in]{geometry}
\usepackage{amsmath}
\usepackage{amssymb}
\usepackage{enumitem}
\usepackage[hidelinks]{hyperref}
\usepackage{xcolor}

\title{CSCE 355: JFLAP3 --- Building DFAs}
\author{}
\date{}

\newcommand{\fn}[1]{\texttt{\detokenize{#1}}}

\begin{document}
\maketitle

% ------------------------------------------------------------------
% NOTE TO STUDENTS: every machine you submit MUST be a DFA. A machine
% that is not deterministic (or has lambda transitions, or reads more
% than one symbol per transition) receives NO CREDIT for that problem.
% See "What counts as a DFA" below for how to check in JFLAP.
% ------------------------------------------------------------------

\section*{Overview}

In this assignment you will build seven deterministic finite automata (DFAs) in
JFLAP. Each is checked automatically. The assignment is worth \textbf{40 points}:
Problems 1--6 are worth 5 points each and Problem 7 is worth 10 points.

\subsection*{Getting started}

\begin{itemize}
  \item Read the JFLAP tutorial: \url{https://jflap.org/tutorial/}.
  \item JFLAP is a free \texttt{.jar} file from \url{https://www.jflap.org}. You need
        Java 8 or higher. Start it with \texttt{java -jar JFLAP7\_1.jar} (use the name of
        the file you downloaded).
  \item Unzip this assignment folder. (It is a good idea to append your username to the
        folder name, e.g.\ \fn{JFLAP3_jsmith}.)
  \item In JFLAP use \textbf{File $\to$ Open} to open one of the provided \texttt{.jff}
        files. Each one is an empty \emph{Finite Automaton}. Build your DFA, then
        \textbf{File $\to$ Save} it \emph{under the same file name}. Do not change the
        machine type.
  \item JFLAP shows the empty string as $\lambda$. Just leave the empty string as
        $\lambda$; do not change it to $\epsilon$.
\end{itemize}

\begin{center}
\begin{tabular}{|c|l|l|c|}
\hline
\textbf{\#} & \textbf{Your DFA (edit this file)} & \textbf{Test strings (do not edit)} & \textbf{Points}\\
\hline
1 & \fn{problem1_01then10.jff}     & \fn{problem1_01then10_tests.txt}     & 5 \\
2 & \fn{problem2_odd1s_0to1.jff}   & \fn{problem2_odd1s_0to1_tests.txt}   & 5 \\
3 & \fn{problem3_no010.jff}        & \fn{problem3_no010_tests.txt}        & 5 \\
4 & \fn{problem4_div7.jff}         & \fn{problem4_div7_tests.txt}         & 5 \\
5 & \fn{problem5_parens.jff}       & \fn{problem5_parens_tests.txt}       & 5 \\
6 & \fn{problem6_01_011.jff}       & \fn{problem6_01_011_tests.txt}       & 5 \\
7 & \fn{problem7_01_110_011.jff}   & \fn{problem7_01_110_011_tests.txt}   & 10 \\
\hline
\end{tabular}
\end{center}

\section*{What counts as a DFA}

\textbf{Every machine you submit must be a DFA.} If a machine is not a DFA it receives
\textbf{zero credit} for that problem, no matter which strings it happens to get right. This
is checked automatically, before any test string is run. A machine is a DFA when:

\begin{itemize}
  \item each state has \textbf{exactly one} outgoing transition on \emph{each} symbol of the
        problem's alphabet: never two (two arrows leaving a state on the same symbol, even to
        different states, is nondeterminism) and never zero (a missing arrow is not allowed);
  \item there are \textbf{no $\lambda$ (empty-string) transitions};
  \item every transition reads \textbf{exactly one symbol}, and that symbol belongs to the
        problem's alphabet (a single arrow may carry both symbols, e.g.\ labeled
        \texttt{0,1}, but a label such as \texttt{01} is \emph{not} allowed);
  \item there is exactly one initial state.
\end{itemize}

\textbf{You must draw the trap (dead) state.} JFLAP lets you leave transitions out and simply
rejects a string that gets stuck, but for this assignment that is \emph{not} a DFA: every state
needs an arrow for every symbol. Wherever a string should be rejected for good, send the missing
arrows to a non-final trap state that loops to itself on every symbol. A DFA missing any
transition receives zero credit for that problem, and the report lists the states and symbols
that are missing.

\subsection*{How to check in JFLAP}

\textbf{Test $\to$ Highlight Nondeterminism.} With your machine open in the editor, choose
this menu item. JFLAP highlights every state that has nondeterministic behavior. If nothing
is highlighted your machine passes this check; if a state is highlighted, look at its
outgoing arrows and fix the conflict. Also make sure you did not leave any $\lambda$
transition behind. When you single-step a string with \textbf{Input $\to$ Step}, a DFA has
exactly one current state at every point. JFLAP does \emph{not} warn you about missing
transitions, so also check by eye that every state has an outgoing arrow for each symbol
(for the alphabet $\{0,1\}$: two arrows per state; the trap state counts as a state). Running
\texttt{python3 test\_dfa.py} (below) checks this too.

\section*{Problems}

\subsection*{Problem 1 (5 points)}
Construct a DFA that recognizes any number of \texttt{01}, followed by \texttt{10}s, without
overlap (e.g.\ \texttt{010} is not in the language recognized by the DFA). The alphabet is
$\{0,1\}$.
\begin{itemize}
  \item The test file has 8 strings: the first 3 should accept, the last 5 should reject.
  \item The empty string should accept (see ``Testing the empty string'' below).
\end{itemize}

\subsection*{Problem 2 (5 points)}
Construct a DFA in 8 or fewer states (I am fine if you delete unreachable states) that is the
product construction of two DFAs: one that recognizes an odd number of \texttt{1}s, and one
that recognizes a binary number with a leading \texttt{0} and ending \texttt{1}. A string is
accepted only if it satisfies \emph{both} conditions. The alphabet is $\{0,1\}$.
\begin{itemize}
  \item The test file has 9 strings: the first 4 should accept, the last 5 should reject.
\end{itemize}

\subsection*{Problem 3 (5 points)}
Construct a DFA that recognizes strings without \texttt{010} as a substring, or that are just
\texttt{01} repeated (\texttt{01}, \texttt{0101}, \texttt{010101}, \ldots). The empty string
should be accepted. The alphabet is $\{0,1\}$.
\begin{itemize}
  \item The test file has 8 strings: the first 4 should accept, the last 4 should reject.
\end{itemize}

\subsection*{Problem 4 (5 points)}
Construct a DFA to recognize unsigned binary strings evenly divisible by 7. The alphabet is
$\{0,1\}$.
\begin{itemize}
  \item You may accept the empty string (it will not be tested).
  \item Leading zeros have no meaning (this is not octal notation): \texttt{0101} and
        \texttt{0000101} are both the number 5.
  \item The test file has 23 strings: the last five should accept, all previous ones reject.
\end{itemize}

\subsection*{Problem 5 (5 points)}
Construct a DFA to recognize, for the alphabet of just parentheses $\Sigma=\{\texttt{(},\texttt{)}\}$,
those strings that have balanced, proper parentheses pairs, with no more than nesting depth 4.
For example, you do not have to recognize \texttt{((((()))))} as being in the language (you
should reject it), but you do have to recognize \texttt{(((())))} or even
\texttt{((())((()))((())))}.
\begin{itemize}
  \item You may accept the empty string (it will not be tested).
  \item The test file has 11 strings: the first 5 should accept, the last 6 should reject.
  \item Once you are done, ask yourself what is interesting about this, do you think? (Not graded.)
\end{itemize}

\subsection*{Problem 6 (5 points)}
Construct a DFA recognizing the concatenation of any number of \texttt{01} and/or \texttt{011},
zero or more copies of each in any order. The alphabet is $\{0,1\}$.
\begin{itemize}
  \item The test file has 8 strings: the first 4 should accept, the last 4 should reject.
  \item The empty string should accept.
\end{itemize}

\subsection*{Problem 7 (10 points)}
Construct a DFA recognizing the concatenation of any number of \texttt{01}, \texttt{110},
and/or \texttt{011}, zero or more copies of each in any order (the copies may be interleaved;
look at the tests). The alphabet is $\{0,1\}$.
\begin{itemize}
  \item The test file has 117 strings: all but the last ten should accept, and the last ten
        should reject.
  \item The empty string should accept.
\end{itemize}

\section*{Testing your DFAs}

\subsection*{In JFLAP}
\begin{enumerate}
  \item Open your DFA (it needs an initial state to be tested).
  \item Choose \textbf{Input $\to$ Multiple Run}.
  \item Click \textbf{Load Inputs} (at the bottom) and select the problem's test file.
  \item Click \textbf{Run Inputs} and compare each row with what the problem says should
        accept and reject.
\end{enumerate}

\subsubsection*{Testing the empty string}
The test files do not contain the empty string. In Multiple Run, put the cursor on the last
row and press Enter to add a blank line (an empty input), then run again. The empty string
should accept in Problems 1, 3, 6, and 7 (and should reject in Problem 2).

\subsection*{From the command line}
This folder contains \texttt{test\_dfa.py}, the same script that is used to grade you (Python 3
is required). From inside this folder run
\begin{center}
\texttt{python3 test\_dfa.py}
\end{center}
to check all seven problems, or \texttt{python3 test\_dfa.py 3} to check only Problem 3. The
report tells you, for each problem: whether the file is really a DFA, which provided test
strings are wrong, and your score. Run it as often as you like.

\section*{Grading}

Each problem earns its points only if \emph{all} of the following hold; there is no partial
credit within a problem.
\begin{itemize}
  \item The machine is a DFA (see above).
  \item Every string in the problem's test file gets the expected result.
  \item Your DFA recognizes \emph{exactly} the described language. The provided strings are
        only examples: the grader compares your DFA with the described language on all
        strings, so a DFA that passes the samples but is wrong somewhere else earns no
        credit. When this happens the report shows a shortest string on which your DFA is
        wrong.
  \item Problem 2 only: at most 8 states are reachable from the initial state (the trap state
        counts if it is reachable).
\end{itemize}

\textbf{Do not edit the test files or \texttt{test\_dfa.py}.} They are replaced with the
original versions before grading, so changing them can only make your own testing misleading.

\section*{Submission}

\begin{enumerate}
  \item Open \textbf{YOURUSERNAME.txt} and put your university username as the \emph{first
        line} of the file (only your username, nothing else on that line). This is how your
        submission is identified.
  \item Make sure each of your seven DFAs is saved in its provided \texttt{.jff} file. Do not
        rename any of the provided files.
  \item Run \texttt{python3 test\_dfa.py} one last time to make sure everything is as you
        expect.
  \item Submit this entire folder (or a zip of it) as instructed.
\end{enumerate}

\end{document}
